\section{Pembahasan}
\subsection{Rekonstruksi Komponen Instrumen dari Proyek FLP}
Untuk membandingkan instrumen, dibutuhkan track dengan not dan waktu
sesuai sumber. Audit awal menemukan satu FLP, dua MP3 dan 12 sampel
\texttt{Classical Acoustic}; belum ada WAV track instrumen. ZIP asli
tidak ditemukan, tetapi paket yang telah diekstrak cukup untuk membaca
proyek. Sampel Classical Acoustic memakai format tag \texttt{0x674f}:
payload Ogg Vorbis di dalam wadah RIFF/WAVE. Sampel berhasil didekode,
tetapi tidak langsung dianggap sebagai track lagu penuh.

PyFLP membaca channel dan not \cite{pyflp}. Proyek berformat FL Studio
25.1.6.4997, tempo 156 BPM, PPQ 96 dan birama 4/4. Tujuh pattern yang
ditemukan adalah Brasses 2, Brasses, Piano, Guitar, Guitar Melody,
Bass dan Tuba. Pembacaan Playlist bawaan PyFLP 2.2.1 tidak sesuai
versi ini, sehingga event Playlist dibaca langsung sebagai 12 record
berukuran 80 byte. Perubahan format juga dicatat pada proyek PyFLP
\cite{playlist}. Indeks dan batas klip diperiksa terhadap data not
serta channel, bukan ditebak dari nama.
Penanda record bernilai 20480 dan indeks pattern/channel diperiksa
sebelum penempatan not. Offset pattern dinyatakan dalam tick;
nilai $-1$ diperlakukan sebagai penanda tanpa pemotongan.

Posisi not memperhitungkan offset pemotongan:
\begin{equation}
 p_{\mathrm{global}}=p_{\mathrm{klip}}+p_{\mathrm{not}}-p_{\mathrm{offset}},
 \qquad t=\frac{p_{\mathrm{global}}}{\mathrm{PPQ}}\frac{60}{\mathrm{BPM}}.
 \label{eq:timing}
\end{equation}
Contohnya, Guitar Melody ditempatkan pada tick 4224 dan dipotong mulai
tick sumber 4224. Not pertama pada tick sumber 4272 harus tetap mulai
di tick global 4272, bukan 8496. Tiap pattern ditempatkan sekali pada
Playlist ini; pengulangan frasa sudah tertulis pada notnya. Tidak
ditemukan perubahan tempo pada controller pattern, sehingga digunakan
tempo tetap yang tersimpan.

Not diekspor menjadi MIDI tiap kelompok dan MIDI multitrack, lalu
dirender melalui FluidSynth 2.6.1 \cite{fluidsynth}. GeneralUser GS
2.0.3 karya S.~Christian Collins menjadi SoundFont pengganti;
lisensi dan sumber disertakan \cite{soundfont}. Pemetaan
Tabel~\ref{tab:preset} mengikuti keluarga instrumen. Pitch dan timing
dipertahankan, tetapi timbre, envelope, artikulasi dan efek asli
tidak dianggap sama.
\begin{table}[H]
\centering\small
\caption{Preset pengganti; program GM dihitung dari nol.}\label{tab:preset}
\begin{tabular}{llll}
\toprule
Kelompok & Channel FLP & Preset asli & SoundFont \\
\midrule
Piano & 3, 4, 5 & Close Grand/lite & Grand Piano (0) \\
Guitar & 2, 6, 7 & Ace Nylon, Classical Acoustic & Nylon Guitar (24) \\
Bass & 9 & Bass & Acoustic Bass (32) \\
Brass & 0 & Full Brass Sustain & Brass Section (61) \\
Tuba & 10 & Tubas Staccato & Tuba (58) \\
\bottomrule
\end{tabular}
\end{table}
Channel Horns Sustain yang tidak memuat not pada pattern terpasang
tidak diberi not baru. Velocity FL maksimum 128 dipetakan proporsional
ke MIDI 0--127; nilai mentah tetap disimpan. Semua render memakai gain
0,2, volume MIDI CC7 sebesar 100, pan CC10 sebesar 64, serta reverb
dan chorus nonaktif. Nomor pitch mentah dipakai sebagai nomor MIDI;
perbedaan penamaan oktaf antarperangkat tidak dijadikan alasan
transposisi. Root kanal bernilai 60, pitch shift kanal/master nol,
dan fine pitch not netral pada sumber yang diperiksa.
Volume mixer dan automation FL Studio tidak direplikasi. Rekonstruksi
menjaga struktur musikal yang terbaca sambil menyatakan bagian
pembentukan suara yang berubah.

\subsection{Validasi Waktu, Durasi, dan Integritas Sinyal}
Sebelum menafsirkan spektrum, perlu dipastikan bahwa data tidak kosong
atau salah tempat. MIDI dibaca ulang dan dibandingkan dengan daftar
not pada pitch, velocity hasil konversi, kanal, serta tick mulai dan
selesai. Pemeriksaan mencakup 2219 not. WAV diperiksa agar amplitudo
tidak nol, sampel valid dan tidak terjadi clipping.
\begin{table}[H]
\centering\small
\caption{Data not dan waktu aktivitasnya. Titik pada tabel numerik
digunakan sebagai pemisah desimal.}\label{tab:dataset}
\begin{tabular}{lrrrr}
\toprule
Instrumen & Not & MIDI min--maks & Awal (s) & Akhir not (s) \\
\midrule
Bass & 187 & 26--39 & 0.962 & 69.034 \\
Piano & 1102 & 38--84 & 1.538 & 69.223 \\
Guitar & 735 & 47--84 & 1.731 & 69.171 \\
Brass & 8 & 55--74 & 0.000 & 32.308 \\
Tuba & 187 & 38--51 & 0.962 & 69.034 \\
\bottomrule
\end{tabular}

\end{table}
Timeline berakhir pada tick 17280, sekitar 69,230769 s pada 156 BPM.
Kelima WAV memiliki 3.141.274 sampel, stereo 44.100 Hz, dan durasi
71,230703 s setelah dua detik tambahan untuk peluruhan. Pembulatan
tempo MIDI menghasilkan selisih kurang dari 70 mikrodetik pada akhir
timeline. Onset audio terukur sekitar 0,27--11,21 ms setelah note-on,
bergantung preset dan ambang deteksi. Selisih dapat berasal dari
attack dan blok synthesizer, bukan perubahan posisi not pada MIDI.

DRUMS dan VOCALS penuh masing-masing berdurasi 160,77 s. Offset klip
audio FLP mengisyaratkan cuplikan sekitar 0,256--69,487 s, tetapi
stretching dan render audio asli tidak diverifikasi. Kedua MP3 penuh
tidak digabung dengan WAV baru. \texttt{All\_Instruments.wav} hanya
menjumlahkan kelima WAV yang memiliki asal waktu dan panjang sama.
Validasi ini memastikan dasar perbandingan timing, tetapi tidak
membuktikan kesamaan bunyi terhadap video YouTube.

\subsection{Analisis Sinyal dalam Domain Waktu}
Analisis waktu ingin mengetahui kapan komponen aktif dan bagaimana
amplitudonya berubah. Osilasi cepat sulit dibaca pada skala seluruh
lagu, sehingga waveform diringkas memakai minimum--maksimum setiap
10 ms dan dibandingkan dengan RMS per 100 ms. Stereo dirata-ratakan
menjadi mono. Rerata global mono dikurangi untuk FFT dan STFT;
tidak ada normalisasi puncak per instrumen.
\begin{figure}[H]
\centering\includegraphics[width=\linewidth]{01_waveform_rms_perbandingan.png}
\caption{Kolom kiri: amplitudo terhadap waktu. Kolom kanan: RMS lokal.
Skala antarbaris sama. Grafik bukan ukuran keras suara yang telah
dikalibrasi.}\label{fig:waktu}
\end{figure}
Pada Gambar~\ref{fig:waktu}, X adalah waktu dari awal Playlist.
Kurva semakin jauh dari nol berarti amplitudo sesaat lebih besar;
RMS semakin tinggi berarti amplitudo efektif potongan lebih besar.
Bass, piano, gitar dan tuba berulang sepanjang sebagian besar
aransemen. Brass terutama muncul pada awal dan sekitar 29,23--32,31 s,
sesuai dua akor pada not sumber. Pola ini menjelaskan mengapa rerata
seluruh rekaman perlu dibedakan dari rerata saat aktif.
\begin{table}[H]
\centering
\caption{RMS global dan RMS pada blok aktif.}\label{tab:rms}
\begin{tabular}{lrrr}
\toprule
Instrumen & RMS seluruh & RMS aktif & Aktif (s) \\
\midrule
Bass & 0.009556 & 0.009744 & 68.50 \\
Piano & 0.014536 & 0.014877 & 68.00 \\
Guitar & 0.010249 & 0.010552 & 67.20 \\
Brass & 0.006334 & 0.023672 & 5.10 \\
Tuba & 0.010588 & 0.011214 & 63.50 \\
\bottomrule
\end{tabular}

\end{table}
Blok aktif ditentukan dari RMS mono sebelum pengurangan rerata yang
sekurang-kurangnya 1\% dari RMS blok tertinggi instrumen tersebut
(ambang relatif $-40$ dB). Ini definisi operasional, bukan deteksi
persepsi atau durasi not persis. Pemeriksaan sebelum pengurangan
rerata mencegah bagian nol dianggap aktif karena offset baru.

Brass memiliki RMS seluruh rekaman sekitar 0,00633, tetapi RMS aktif
sekitar 0,02367. Nilai global kecil terutama dipengaruhi banyaknya
bagian diam. Pada piano, nilainya lebih dekat: sekitar 0,01454 dan
0,01488, karena aktivitas lebih merata. Secara matematis,
rerata kuadrat pada~\eqref{eq:rms} memasukkan bagian senyap ketika
seluruh timeline dihitung. Karena itu, hasil tidak cukup untuk
menyatakan ``brass lebih pelan daripada piano''.

Untuk memeriksa awal bunyi, Gambar~\ref{fig:attack} memakai 180 ms
sejak note-on pertama tiap kelompok. Y adalah RMS 5 ms, dengan skala
berbeda antarbaris agar bentuk awal terbaca. Garis putus adalah
note-off, bukan saat audio harus langsung nol.
\begin{figure}[H]
\centering\includegraphics[width=.92\linewidth]{04_attack_awal.png}
\caption{RMS sejak not pertama. Note-off piano, bass dan tuba sekitar
92 ms; gitar sekitar 132 ms. Akor brass masih ditahan sepanjang
tampilan.}\label{fig:attack}
\end{figure}
Piano sudah memiliki RMS nyata pada belasan milidetik awal, sedangkan
bass, gitar dan tuba pada contoh ini meningkat selama puluhan
milidetik awal. Sesudah note-off sekitar 92 ms, tuba mendekati nol
lebih cepat daripada sisa amplitudo piano dan bass pada rentang
tampilan. Brass mempertahankan amplitudo karena akornya belum dilepas.
Ini contoh perbedaan awal dan peluruhan pada render, bukan sifat
mutlak setiap instrumen.

RMS 5 ms masih dipengaruhi fase osilasi, terutama nada rendah dengan
periode melebihi 5 ms. Riak dan waktu puncak kecil tidak ditafsirkan
sebagai konstanta attack murni. Grafik menghubungkan envelope dengan
pembobotan energi lokal, tetapi not, velocity dan jumlah suara
bersamaan pada kelima contoh memang berbeda.

\subsection{Analisis Fourier Global dan Distribusi Energi}
Analisis ini ingin membandingkan bagian spektrum yang berkontribusi
kuat. FFT diterapkan pada seluruh sinyal mono setelah pengurangan
rerata. Tanpa jendela tambahan, batas rekaman bertindak sebagai
jendela persegi panjang. Jarak bin kelima WAV sekitar 0,014039 Hz.
Puncak dipilih dari maksimum lokal pada 20--10.000 Hz dengan jarak
minimal 10 Hz agar bukan hanya bin tetangga dari satu gugus.
\begin{figure}[H]
\centering\includegraphics[width=\linewidth]{02_fft_perbandingan.png}
\caption{X: frekuensi; Y: kerapatan daya dalam dB. Tampilan merangkum
maksimum per 0,5 Hz; hitungan memakai bin asli. Skala panel sama.}
\label{fig:fft}
\end{figure}
Kurva tinggi pada Gambar~\ref{fig:fft} berarti kontribusi spektral
kuat. Bass terkonsentrasi dekat sisi kiri, sedangkan brass memiliki
puncak kuat lebih jauh ke frekuensi tinggi. Piano dan gitar mempunyai
beberapa gugus karena memainkan banyak nada beserta harmoniknya.
\begin{table}[H]
\centering\small
\caption{Ringkasan WAV baru. Puncak adalah puncak terpilih terkuat,
bukan penetapan fundamental.}\label{tab:spektrum}
\begin{tabular}{lrrrr}
\toprule
Instrumen & Puncak (Hz) & Sentroid (Hz) & $<250$ Hz (\%) & $\geq1000$ Hz (\%) \\
\midrule
Bass & 49.40 & 91.54 & 97.61 & 0.05 \\
Piano & 312.00 & 590.64 & 3.40 & 8.54 \\
Guitar & 195.00 & 575.14 & 17.54 & 9.43 \\
Brass & 1174.61 & 1323.02 & 3.09 & 55.70 \\
Tuba & 117.00 & 309.67 & 41.04 & 0.52 \\
\bottomrule
\end{tabular}

\end{table}
Bass memiliki puncak sekitar 49,40 Hz dan 97,61\% energi di bawah
250 Hz. Tuba memiliki puncak sekitar 117,00 Hz, tetapi hanya 41,04\%
energi pada pita tersebut. Data MIDI menunjukkan 187 not bass dan
tuba memiliki 187 onset yang sama, dengan pitch tuba satu oktaf
lebih tinggi. Sebanyak 20 note-off tuba lebih lambat 2 atau 4 tick,
sehingga waktu akhir not tidak seluruhnya identik.
Dalam rumus nada ideal $\nu(p)=440\,2^{(p-69)/12}$, kenaikan 12
semiton menggandakan frekuensi. Namun, spektrum tuba bukan sekadar
salinan bass yang digeser karena preset dan kekuatan harmonik berbeda.

Brass memiliki puncak sekitar 1174,61 Hz, meskipun not tertinggi MIDI
74 berfrekuensi ideal sekitar 587,33 Hz. Puncak tersebut dekat
harmonik kedua nada itu. Akor memuat beberapa not, sehingga harmonik
juga dapat bertumpuk. Inilah contoh puncak terbesar yang tidak langsung
menjadi fundamental. Pada bass, puncak sekitar 49,40 dan 98,81 Hz
mendukung hubungan harmonik mendekati 1:2, tetapi spektrum seluruhnya
tetap gabungan banyak kejadian.

Sentroid dan fraksi pita merangkum bentuk distribusi daya:
\begin{equation}
 C=\frac{\sum_m\nu_mP_m}{\sum_mP_m},\qquad
 \eta_B=\frac{\sum_{\nu_m\in B}P_m}{\sum_mP_m}.
 \label{eq:fitur}
\end{equation}
Sentroid lebih tinggi berarti pusat daya bergeser ke frekuensi lebih
tinggi, bukan amplitudo lebih besar. Bass mempunyai sentroid sekitar
91,54 Hz, brass sekitar 1323,02 Hz. Secara spektral, brass lebih kaya
komponen tinggi dan dapat dianalogikan relatif lebih terang; ini
bukan hasil uji pendengaran subjektif.
\begin{figure}[H]
\centering\includegraphics[width=.94\linewidth]{03_pita_energi.png}
\caption{X: kelompok instrumen. Tinggi segmen warna: persentase energi
pada pita di legenda. Tiap batang berjumlah 100\%, sehingga menekankan
komposisi, bukan besar amplitudo keseluruhan.}\label{fig:pita}
\end{figure}
Sekitar 55,70\% energi brass berada di atas 1000 Hz, sedangkan bass
hanya sekitar 0,05\%. Piano dan gitar masing-masing 8,54\% dan 9,43\%.
Pembagian mengikuti penjumlahan daya pada~\eqref{eq:fitur}; batas
pita bukan kategori instrumen universal. Perbandingan dipengaruhi
register, partitur, velocity, lama aktif dan preset, sehingga tidak
mengisolasi timbre sebagai satu-satunya penyebab.

\subsection{Validasi Energi Menggunakan Identitas Parseval}
Pemeriksaan ini ingin mengetahui apakah normalisasi dan pembobotan
spektrum benar. Energi dihitung langsung dari sampel dan koefisien
FFT dengan~\eqref{eq:parseval}.
\begin{table}[H]
\centering\small
\caption{Energi dalam amplitudo digital kuadrat detik.}\label{tab:parseval}
\begin{tabular}{lrrr}
\toprule
Instrumen & $E_t$ & $E_F$ & Galat relatif \\
\midrule
Bass & 0.006504213 & 0.006504213 & 5.33e-16 \\
Piano & 0.015049983 & 0.015049983 & 8.07e-16 \\
Guitar & 0.007482171 & 0.007482171 & 1.28e-15 \\
Brass & 0.002857970 & 0.002857970 & 6.22e-15 \\
Tuba & 0.007986120 & 0.007986120 & 1.95e-15 \\
\bottomrule
\end{tabular}

\end{table}
Galat relatif semua instrumen di bawah $7\times10^{-15}$, setingkat
pembulatan numerik. Validasi STFT pada tiga bingkai representatif
per jendela juga menghasilkan galat di bawah $4\times10^{-16}$.
Bagi interpretasi audio, ini mendukung bahwa energi tidak berubah
karena salah faktor FFT. Parseval tidak memvalidasi kesamaan SoundFont
terhadap preset asli atau mengidentifikasi instrumen: teorema tersebut
memastikan konsistensi dua representasi sinyal yang sama.
